\chapter{Conclusion}\label{ch:conclusion}

\paragraph{Summary.} The project implemented a complete offline multiscale
pipeline for rigid-particle-reinforced PDMS in Python. \Cref{tab:status} separates
what was achieved, what was implemented without fully meeting its objective, and
what is proposed but not implemented.

{\small
\begin{xltabular}{\textwidth}{@{}p{2.6cm}X@{}}
\caption{Status of the project components.}\label{tab:status}\\
\toprule
Status & Components\\
\midrule
\endfirsthead
\caption[]{\emph{(continued)}}\\
\toprule
Status & Components\\
\midrule
\endhead
\bottomrule
\endfoot
Achieved (implemented and verified) &
  Periodic finite-strain RVE with mixed near-incompressible matrix and exactly
  rigid, freely rotating particle; reaction-based Hill--Mandel homogenization;
  exact Schur-condensed homogenized tangent; reproducible, leak-free dataset
  generation; polyconvex, volume-fraction-conditioned neural energy with objective
  normalization; PyTorch--NGSolve bridge with energy, stress and full $9\times9$
  tangent; 26 passing verification tests (\cref{tab:vmatrix}); RVE moduli within
  $1.1\,\%$ of Mori--Tanaka; median speedup of $5.8\times10^4$ for repeated RVE
  evaluations in the tested setting (\cref{ch:performance}).\\
Partially achieved &
  Surrogate accuracy: accurate in energy and pressure, but not in deviatoric stress
  (\cref{sec:accuracy}); shear fine-tuning reduced but did not remove the error.
  Macroscopic application: $P_1$ reaches $5\,\%$ stretch but is locked, $P_2$ reaches
  $1\,\%$ on two of three meshes (\cref{sec:macro}). No experimental validation.\\
Proposed (not implemented) &
  Volumetric--isochoric surrogate with balanced loss and targeted sampling
  (\cref{sec:proposed}); mixed $u$--$p$ macroscopic formulation (\cref{sec:macro}).\\
\end{xltabular}
}

The speedup and the accuracy are separate results: the speedup shows that the
surrogate is cheap to evaluate, not that its constitutive response is correct.

\paragraph{Main findings.}
\begin{enumerate}
  \item The surrogate is accurate in pressure (median error $0.03\,\%$) but not in
        deviatoric stress (median $76$--$83\,\%$ outside expansion states). Its small-strain shear modulus is
        $46$--$70\,\%$ too high. The cause is quantified: with $\kappa/\mu=10^4$,
        volumetric terms dominate the training loss, and deleting the whole
        deviatoric response changes the stress loss by only $1.5\,\%$. The headline
        training-scale NRMSE of $\sim10^{-3}$ is correct but blind to this error.
        Along five loading paths the four paths with a volume change agree with the
        direct RVE to better than $1\,\%$ in stress, while the simple-shear error is
        constant from $2.5$ to $10\,\%$ amplitude: a wrong modulus, not a missing
        nonlinearity (\cref{sec:paths}).
  \item Shear fine-tuning reduced the shear-holdout error from $59\,\%$ to $40\,\%$ by
        lowering $\mu_{\mathrm{eff}}$. It did not change the general deviatoric
        error, consistent with the diagnosis.
  \item The $P_1$ macroscopic solution to $5\,\%$ stretch is volumetrically locked
        (energy 25--89 times the $P_2$ value). The $P_2$ solution is plausible but
        hard to drive beyond $1\,\%$ because of the $\kappa/\mu$ conditioning of a
        displacement-only formulation.
  \item The computational bottlenecks are the offline RVE solves (62 single-core
        hours, poorly balanced across workers) and, online, the Newton iteration
        count of the stiff macroscopic problem (2005--3084 iterations for $P_1$); the
        AD tangent dominates only the cost of each iteration (\cref{ch:performance}).
  \item Testing caught genuine defects that accuracy metrics could not reveal. The
        Phase~1 network is not frame-indifferent (three negative stiffness
        eigenvalues, V0d) and its shear modulus is $12\,\%$ too low (V8), hidden behind a
        $3\,\%$ stress error. The Phase~5 surrogate failed objectivity (S-V1).
\end{enumerate}

\paragraph{Outlook.} The next steps follow directly from these findings:
(i) the volumetric--isochoric energy \cref{eq:split-proposed} with the balanced
loss \cref{eq:loss-proposed}; (ii) sampling restricted to
physically relevant $J$ with dense isochoric coverage; (iii) a mixed $u$--$p$ or
$\bar{\bm F}$ macroscopic formulation, re-using the RVE's Taylor--Hood
implementation, which requires (i) and its own verification (patch test, mesh
refinement, Newton convergence, comparison with direct RVE solutions); (iv) an analytical or vectorized tangent, which removes
about 80\,\% of the online cost; and (v) several particles per cell, a random
microstructure, and finally comparison with experimental PDMS--glass data. None
of these steps requires changes to the verified RVE, homogenization or bridge
infrastructure, which is the main durable result of this project.
